% !TEX root = ../Main.tex
\chapter{蛙泳}

\section{出发与转身后动作}

\state{出发 / 转身后允许的动作}{
    \begin{itemize}
        \item 允许身体\textbf{没入水中}；
        \item 可做\textbf{一次手臂充分向后划至腿部}的动作；
        \item 在第一次手臂划水动作过程中，允许\textbf{打一次蝶泳腿}接\textbf{蛙泳蹬腿}动作。
    \end{itemize}
}

\section{整体姿势}

\state{俯卧与动作周期}{
    \begin{itemize}
        \item 从出发和每次转身后的\textbf{第一次手臂动作}开始，\textbf{身体应保持俯卧}；
        \item 任何时候\textbf{都不允许身体呈仰卧姿势}；
        \item 在出发后的整个游程中，动作周期必须是\textbf{"一次划臂 $+$ 一次蹬腿"}的顺序完成。
    \end{itemize}
}

\section{手臂动作}

\state{两臂的对称要求}{
    \begin{itemize}
        \item 两臂所有动作\textbf{应同时并在同一水平面上进行}，\textbf{不得有交替动作}；
        \item 两手应\textbf{同时在水面、水下或水上由胸前伸出}；
        \item \textbf{肘部不得露出水面}；例外情况：\textbf{转身前最后一次划水、转身过程中、抵达终点前最后一次划水}；
        \item 两手应\textbf{在水面或水下向后划水}；
        \item 除\textbf{出发和每次转身后的第一次划水}外，\textbf{两手向后划水不得超过臀线}。
    \end{itemize}
}

\section{头部规则}

\state{头部必须露出水面}{
    \begin{itemize}
        \item 在\textbf{每个完整动作周期内}，运动员\textbf{头的某一部分必须露出水面}；
        \item 出发和每次转身后\textbf{第二次划臂}至\textbf{最宽点两手向内划水前}，头部必须\textbf{露出水面}。
    \end{itemize}
}

\section{腿部动作}

\state{两腿的对称与外翻}{
    \begin{itemize}
        \item 两腿所有动作应\textbf{同时并在同一水平面上进行}，\textbf{不得有交替动作}；
        \item 在蹬腿过程中，两脚必须做\textbf{外翻动作}；
        \item \textbf{不允许剪夹、上下交替打水或向下的蝶泳打水动作}（出发 / 转身后允许的那一次蝶泳腿除外）；
        \item 只要不做\textbf{向下的蝶泳打腿}，允许两脚\textbf{露出水面}。
    \end{itemize}
}

\section{转身与终点}

\state{双手触壁}{
    \begin{itemize}
        \item 在每次\textbf{转身}和\textbf{到达终点}时，两手应\textbf{在水面、水上或水下同时触壁}；
        \item 触壁前\textbf{最后一次划水结束后}，头可以没入水中；
        \item 但在触壁前\textbf{最后一个完整或不完整的动作周期}中，头的某一部分\textbf{应露出水面}。
    \end{itemize}
}

\rmkb{
    \textbf{易考点}：
    \begin{itemize}
        \item 动作周期：\textbf{一次划臂} $+$ \textbf{一次蹬腿}；
        \item 两臂两腿\textbf{同时、同水平面、对称}，不得交替；
        \item \textbf{蹬腿要外翻}；不允许剪夹 / 上下交替打水 / 向下蝶泳打腿；
        \item 出发 / 转身第一次划水\textbf{允许一次蝶泳腿} $+$ 蛙泳蹬腿；
        \item \textbf{双手同时触壁}（转身与终点均然）；
        \item 每个完整动作周期内\textbf{头必须露出水面}。
    \end{itemize}
}
